\documentclass[11pt,a4paper,oneside]{report}

% ------------------------------------------------------------
% PACOTES
% ------------------------------------------------------------

% ------------------------------------------------------------
% PACOTES
% ------------------------------------------------------------
\usepackage[utf8]{inputenc}
\usepackage[T1]{fontenc}
\IfFileExists{brazil.ldf}
  {\usepackage[brazil]{babel}}
  {\usepackage[english]{babel}}

\usepackage{lmodern}
\usepackage{microtype}
\usepackage{geometry}
\usepackage{graphicx}
\usepackage{float}
\usepackage{booktabs}
\usepackage{array}
\usepackage{amsmath,amssymb}
\usepackage{xcolor}
\usepackage{enumitem}
\usepackage{caption}
\usepackage{fancyhdr}
\usepackage{titlesec}
\usepackage{hyperref}
\usepackage{tocloft}

% ------------------------------------------------------------
% CONFIGURACOES GERAIS
% ------------------------------------------------------------
\geometry{top=2.1cm,bottom=2.1cm,left=2.3cm,right=2.3cm}
\graphicspath{{figuras/}{./}}
\setlength{\parindent}{1.2cm}
\setlength{\parskip}{0.25em}
\renewcommand{\arraystretch}{1.18}
\raggedbottom

\definecolor{azulpoli}{RGB}{0,70,127}
\definecolor{cinzaclaro}{RGB}{242,244,247}

\titleformat{\chapter}[hang]
  {\Huge\bfseries\color{black}}
  {\thechapter.}{0.6em}{}

\titlespacing*{\chapter}
  {0pt}
  {-20pt}
  {25pt}

\titleformat{\section}
  {\Large\bfseries\color{black}}
  {\thesection.}{0.6em}{}

\titleformat{\subsection}
  {\large\bfseries\color{black}}
  {\thesubsection}{0.6em}{}

\pagestyle{fancy}
\fancyhf{}
\fancyhead[L]{PSI3431 -- Processamento Estatístico de Sinais}
\fancyhead[R]{Identificação de Sistemas com LMS}
\fancyfoot[C]{\thepage}
\setlength{\headheight}{14pt}

\hypersetup{
    colorlinks=true,
    linkcolor=black,
    urlcolor=black,
    citecolor=black,
    pdftitle={Identificação de Sistemas com o Algoritmo LMS},
    pdfauthor={José Victor Lyra de Castro, Caio de Paula Colnago e Arthur Jesus}
}

\captionsetup{font=small,labelfont=bf}
\renewcommand{\cfttoctitlefont}{\Huge\bfseries}
\renewcommand{\cftsecfont}{\bfseries}
\renewcommand{\cftsecpagefont}{\bfseries}
\renewcommand{\cftsubsecfont}{\normalfont}
\renewcommand{\cftsubsecpagefont}{\normalfont}
\renewcommand{\cftchapfont}{\bfseries}
\renewcommand{\cftchappagefont}{\bfseries}
\newcommand{\wtrue}{w^{o}}
\newcommand{\what}{\widehat{w}}

\begin{document}

% ------------------------------------------------------------
% CAPA
% ------------------------------------------------------------
\hypersetup{pageanchor=false}
\begin{titlepage}
    \thispagestyle{empty}
    \begin{center}
        \IfFileExists{poli.png}{
            \includegraphics[width=0.22\textwidth]{poli.png}
        }{
            \vspace*{1.5cm}
        }

        \vspace{0.7cm}

        {\large\textbf{UNIVERSIDADE DE SÃO PAULO}}\\[0.15cm]
        {\large Escola Politécnica}\\[0.15cm]
        {\large PSI3431 -- Processamento Estatístico de Sinais}

        \vfill

        {\LARGE\bfseries Identificação de Sistemas com o Algoritmo LMS}\\[0.45cm]
        {\large}

        \vfill

        \begin{tabular}{ll}
            \textbf{Integrante 1:} & José Victor Lyra de Castro \\
            \textbf{NUSP:} & 14747340 \\[0.4cm]

            \textbf{Integrante 2:} & Caio de Paula Colnago \\
            \textbf{NUSP:} & 12608900 \\[0.4cm]

            \textbf{Integrante 3:} & Artur Mateo Freire de Jesus \\
            \textbf{NUSP:} & 14745525 \\
        \end{tabular}

        \vfill

        \textbf{Código do projeto:}\\[0.15cm]
        \href{https://github.com/Art-Mat/PES}
        {\texttt{github.com/Art-Mat/PES}}

        \vspace{0.8cm}

        {\large São Paulo}\\
        {\large 2026}
    \end{center}
\end{titlepage}
\hypersetup{pageanchor=true}
\pagenumbering{arabic}

% ------------------------------------------------------------
% SUMARIO
% ------------------------------------------------------------
\thispagestyle{fancy}
\tableofcontents
\vfill

\clearpage

% ------------------------------------------------------------
% TOPICO A - PAGINA 3
% ------------------------------------------------------------

\chapter{Parte 1}

\section{Visão geral da solução}

\paragraph{} O código implementa a identificação de um sistema por meio de um
filtro adaptativo LMS (\textit{Least Mean Squares}). A planta é
conhecida e possui dez coeficientes. Para verificar estatisticamente a
convergência do algoritmo, o experimento é repetido em
\(1000\) realizações, sempre utilizando novas sequências pseudoaleatórias
de entrada e de ruído.

Em cada realização, o algoritmo executa as seguintes etapas:

\begin{enumerate}[label=\arabic*., leftmargin=*, itemsep=0.3em]
    \item gera um sinal de entrada branco;
    \item calcula a saída da planta conhecida;
    \item adiciona ruído à saída da planta;
    \item calcula a saída do filtro adaptativo;
    \item determina o erro entre a saída desejada e a saída estimada;
    \item atualiza os coeficientes do filtro por meio do algoritmo LMS;
    \item armazena os erros necessários para estimar as curvas de
    MSE e EMSE.
\end{enumerate}

\subsection{Organização geral dos arquivos}

\paragraph{} A implementação foi dividida em cinco arquivos, de forma a separar os
parâmetros da simulação, o processamento do algoritmo e a geração dos
gráficos. A organização adotada é apresentada na
Tabela~\ref{tab:organizacao-arquivos}.

\begin{table}[H]
    \centering
    \caption{Organização dos arquivos utilizados na implementação.}
    \label{tab:organizacao-arquivos}

    \small
    \setlength{\tabcolsep}{6pt}
    \renewcommand{\arraystretch}{1.25}

    \begin{tabular}{
        @{}
        >{\raggedright\arraybackslash}p{0.25\textwidth}
        >{\raggedright\arraybackslash}p{0.68\textwidth}
        @{}
    }
        \toprule
        \textbf{Arquivo} & \textbf{Responsabilidade} \\
        \midrule

        \texttt{main.m}
        &
        Iniciar a execução geral do projeto.
        \\

        \texttt{PartOneConsts.m}
        &
        Centralizar os parâmetros e as constantes utilizadas na simulação.
        \\

        \texttt{knownPlant.m}
        &
        Executar o processo de identificação da planta conhecida por meio
        do algoritmo LMS.
        \\

        \texttt{PlotMSECurve.m}
        &
        Calcular a média dos erros quadráticos e gerar a curva de
        aprendizado do MSE.
        \\

        \texttt{PlotEMSECurve.m}
        &
        Calcular a média dos erros quadráticos em excesso e gerar a curva
        de EMSE.
        \\

        \bottomrule
    \end{tabular}
\end{table}

Essa separação torna a implementação mais modular e facilita sua
manutenção. Os parâmetros da simulação ficam concentrados em um único
arquivo, enquanto o processamento do algoritmo e as rotinas de
visualização são mantidos em funções independentes. Dessa forma,
alterações nos parâmetros, no algoritmo ou na apresentação dos
resultados podem ser realizadas sem modificar toda a estrutura do
projeto.


\section{Algoritmo, recursão e parâmetros utilizados}

\subsection{Problema de identificação}

A planta considerada é um filtro FIR de dez coeficientes, representado pelo vetor verdadeiro

\begin{equation}
\resizebox{0.96\textwidth}{!}{$
\wtrue=
[-0.5247\; -0.2060\; 0.3324\; -0.2631\; 0.2358\; 0.0304\; 0.3525\; -0.4812\; -0.0485\; -0.2964]^T.
$}
\end{equation}

Embora \(\wtrue\) seja conhecido na simulação, ele é tratado como desconhecido pelo filtro adaptativo.
O objetivo do LMS é obter uma estimativa \(\what(i)\) a partir apenas da entrada e da saída observada.
Para cada instante \(i\), define-se o vetor de entrada

\begin{equation}
    {u}(i)=[u(i)\;u(i-1)\;\cdots\;u(i-9)]^T.
\end{equation}

O sinal \(u(i)\) é branco, gaussiano, de média nula e variância unitária.

Essa escolha é adequada porque o sinal branco excita todas as direções do filtro e, por possuir matriz
de autocorrelação proporcional à identidade, favorece uma convergência uniforme dos coeficientes.

A saída observada da planta é

\begin{equation}
    d(i)={u}^T(i)\wtrue+v(i),
\end{equation}

em que \(v(i)\) é um ruído branco de saída, independente da entrada, com média zero e variância
\(\sigma_v^2=0{,}001\). A saída do filtro adaptativo e o erro de estimação são, respectivamente,

\begin{equation}
    y(i)={u}^T(i)\what(i),
    \qquad
    e(i)=d(i)-y(i).
\end{equation}

\subsection{Critério e recursão do LMS}

O critério a ser minimizado é o erro médio quadrático

\begin{equation}
    J(w)=E[e^2(i)].
\end{equation}

O LMS substitui o gradiente estatístico de \(J(w)\) por uma estimativa instantânea. Com a
convenção adotada no código, a atualização dos coeficientes é

\begin{equation}
    \what(i+1)=\what(i)+\mu\,{u}(i)e(i).
\end{equation}

Essa recursão desloca os coeficientes na direção que reduz o erro quadrático observado. A solução tende
à solução de Wiener \(w^{\star}=R_u^{-1}r_{du}\), sem exigir o cálculo explícito das
correlações.

\clearpage

% ------------------------------------------------------------
% TOPICO A - PAGINA 4
% ------------------------------------------------------------
\subsection{Configuração da simulação}

Os parâmetros utilizados estão resumidos na Tabela~\ref{tab:parametros}.

\begin{table}[H]
    \centering
    \caption{Parâmetros empregados na identificação da planta conhecida.}
    \label{tab:parametros}
    \begin{tabular}{p{5.0cm} p{3.2cm} p{5.2cm}}
        \toprule
        \textbf{Parâmetro} & \textbf{Valor} & \textbf{Finalidade} \\
        \midrule
        Número de coeficientes \(L\) & 10 & Dimensão da planta e do filtro adaptativo \\
        Realizações independentes & 1000 & Aproximação das médias estatísticas \\
        Iterações por realização & 500 & Observação do transitório e do regime permanente \\
        Passo de adaptação \(\mu\) & 0,02 & Compromisso entre velocidade e erro residual \\
        Variância da entrada \(\sigma_u^2\) & 1 & Entrada branca de potência unitária \\
        Variância do ruído \(\sigma_v^2\) & 0,001 & Piso de erro irredutível da observação \\
        Semente aleatória & 0 & Reprodutibilidade numérica \\
        Inicialização & \(\what(0)=0\) & Ausência de conhecimento prévio da planta \\
        \bottomrule
    \end{tabular}
\end{table}

Para cada realização, o vetor de entrada e os coeficientes estimados são zerados. Em cada iteração,
uma nova amostra gaussiana é inserida em \(u(i)\), calcula-se \(d(i)\), obtém-se o erro e aplica-se
a recursão do LMS. As realizações usam sequências independentes de entrada e ruído.

O procedimento pode ser resumido em quatro passos:

\begin{enumerate}[label=\arabic*.,leftmargin=1.0cm,itemsep=0.25em]
    \item gerar \(d(i)=u^T(i)\wtrue+v(i)\);
    \item calcular \(y(i)=u^T(i)\what(i)\) e \(e(i)=d(i)-y(i)\);
    \item armazenar o erro quadrático e o erro em excesso;
    \item atualizar \(\what(i+1)=\what(i)+\mu u(i)e(i)\).
\end{enumerate}

O MSE experimental é obtido pela média dos erros quadráticos nas \(M=1000\) realizações:

\begin{equation}
    \widehat{\mathrm{MSE}}(i)=\frac{1}{M}\sum_{m=1}^{M}e_m^2(i).
\end{equation}

Já o erro em excesso desconsidera diretamente o ruído de saída:

\begin{equation}
    e_{a,m}(i)=u_m^T(i)\bigl(\wtrue-\what_m(i)\bigr),
    \qquad
    \widehat{\mathrm{EMSE}}(i)=\frac{1}{M}\sum_{m=1}^{M}e_{a,m}^2(i).
\end{equation}

Como \(R_u=I\), o maior autovalor é \(\lambda_{\max}=1\). O valor \(\mu=0{,}02\)
é pequeno em relação aos limites usuais de estabilidade, sendo adequado para obter convergência suave
e baixo erro residual, ainda que exija algumas centenas de iterações.

\clearpage

% ------------------------------------------------------------
% TOPICO B - PAGINA 5
% ------------------------------------------------------------
\section{Curva do erro médio quadrático (MSE)}

O MSE inclui tanto o erro decorrente da diferença entre os coeficientes quanto o ruído presente na saída:

\begin{equation}
    \mathrm{MSE}(i)=E[e^2(i)]
    =E[e_a^2(i)]+\sigma_v^2
    =\mathrm{EMSE}(i)+J_{\min},
\end{equation}

em que \(J_{\min}=\sigma_v^2\), pois o ruído é independente da entrada e não pode ser previsto pelo filtro.

\begin{figure}[H]
    \centering
    \IfFileExists{MSE_1.png}{
        \includegraphics[width=0.88\textwidth]{MSE_1.png}
    }{
        \fbox{\parbox[c][7cm][c]{0.85\textwidth}{\centering Inserir o arquivo \texttt{MSE\_1.png}.}}
    }
    \caption{Curva de aprendizado do MSE, calculada sobre 1000 realizações.}
    \label{fig:mse}
\end{figure}

Observa-se um curto crescimento inicial, associado ao preenchimento progressivo da linha de atraso e ao
transitório dos coeficientes inicialmente nulos. Depois desse período, o MSE diminui de forma aproximadamente
monótona e alcança o regime permanente por volta de 300 iterações.

O patamar final situa-se próximo de \(-29{,}5\,\mathrm{dB}\), ligeiramente acima do limite ideal de
\(-30\,\mathrm{dB}\). Esse limite decorre diretamente da variância do ruído:

\begin{equation}
    10\log_{10}(\sigma_v^2)
    =10\log_{10}(0{,}001)
    =-30\,\mathrm{dB}.
\end{equation}

A pequena diferença é esperada: com passo constante, o LMS não permanece exatamente no ponto ótimo,
mas oscila em torno dele, produzindo um erro de desajuste adicional.

\clearpage

% ------------------------------------------------------------
% TOPICO C - PAGINA 6
% ------------------------------------------------------------
\section{Curva do erro médio quadrático em excesso (EMSE)}

O EMSE mede apenas a parcela de erro provocada pelo afastamento entre os coeficientes estimados e os
coeficientes verdadeiros:

\begin{equation}
    \mathrm{EMSE}(i)
    =E\!\left[\left|u^T(i)(\wtrue-\what(i))\right|^2\right].
\end{equation}

Assim, diferentemente do MSE, o EMSE não contém diretamente a potência do ruído de saída.

\begin{figure}[H]
    \centering
    \IfFileExists{EMSE_1.png}{
        \includegraphics[width=0.88\textwidth]{EMSE_1.png}
    }{
        \fbox{\parbox[c][7cm][c]{0.85\textwidth}{\centering Inserir o arquivo \texttt{EMSE\_1.png}.}}
    }
    \caption{Curva do erro médio quadrático em excesso, calculada sobre 1000 realizações.}
    \label{fig:emse}
\end{figure}

A curva apresenta o mesmo transitório inicial observado no MSE, seguido de queda contínua até aproximadamente
300 iterações. O regime permanente ocorre em torno de \(-39{,}5\,\mathrm{dB}\), evidenciando que a parcela
de erro associada ao desajuste dos coeficientes tornou-se aproximadamente dez vezes menor que a potência do ruído.

O EMSE não converge exatamente a zero porque o passo \(\mu\) permanece constante. Mesmo depois de atingir
a vizinhança de \(\wtrue\), cada nova realização instantânea de entrada e ruído continua perturbando a
estimativa. Um passo menor reduziria esse patamar, mas aumentaria o tempo necessário para a convergência.

\clearpage

% ------------------------------------------------------------
% TOPICO D - PAGINA 7
% ------------------------------------------------------------
\section{Análise dos resultados e justificativa das escolhas}

Os resultados confirmam a identificação da planta. Após aproximadamente 300 iterações, MSE e EMSE entram
em regime permanente, indicando que os coeficientes estimados passaram a oscilar em uma pequena vizinhança
dos coeficientes verdadeiros. A média de 1000 realizações reduz a variabilidade das curvas e permite interpretar
os resultados como aproximações das esperanças estatísticas.

Para entrada branca gaussiana de variância unitária, o desajuste aproximado do LMS é

\begin{equation}
    \mathcal{M}\approx\frac{\mu L\sigma_u^2}{2}
    =0{,}10.
\end{equation}

Consequentemente, as previsões de regime permanente são

\begin{align}
    \mathrm{EMSE}_{\infty}
    &\approx \mathcal{M}\,J_{\min}
      =0{,}10\times0{,}001=10^{-4}
      \quad\Rightarrow\quad -40\,\mathrm{dB},\\
    \mathrm{MSE}_{\infty}
    &\approx J_{\min}+\mathrm{EMSE}_{\infty}
      =0{,}0011
      \quad\Rightarrow\quad -29{,}59\,\mathrm{dB}.
\end{align}

\begin{table}[H]
    \centering
    \caption{Comparação entre previsões teóricas e resultados observados.}
    \begin{tabular}{lcc}
        \toprule
        \textbf{Grandeza} & \textbf{Previsão aproximada} & \textbf{Valor observado} \\
        \midrule
        Início do regime permanente & cerca de 300 iterações & cerca de 300 iterações \\
        MSE em regime permanente & \(-29{,}59\,\mathrm{dB}\) & aproximadamente \(-29{,}5\,\mathrm{dB}\) \\
        EMSE em regime permanente & \(-40\,\mathrm{dB}\) & aproximadamente \(-39{,}5\,\mathrm{dB}\) \\
        \bottomrule
    \end{tabular}
\end{table}

A proximidade entre teoria e simulação valida a implementação. O passo \(\mu=0{,}02\) mostrou-se uma escolha
coerente: forneceu convergência estável dentro das 500 iterações e desajuste de aproximadamente 10\%. Um valor
maior aceleraria inicialmente a adaptação, mas elevaria o EMSE de regime permanente e poderia comprometer a
estabilidade. Um valor menor reduziria o erro residual, ao custo de uma convergência mais lenta.

Também foi adequada a utilização de entrada branca, pois \(R_u=I\) reduz a dispersão entre os modos
de convergência. Para entradas coloridas, a maior separação entre os autovalores da matriz de autocorrelação
tenderia a tornar a convergência mais lenta e desigual.

\chapter{Parte 2}

\section{Visão geral da solução}

\paragraph{} Nesta etapa, a resposta ao impulso da planta não é fornecida. A identificação deve ser realizada
exclusivamente a partir dos sinais de entrada e saída contidos nos arquivos \texttt{whites.mat} e
\texttt{correl.mat}. Cada conjunto contém 100 realizações do mesmo sistema, sendo que o primeiro
utiliza entrada branca e o segundo utiliza entrada temporalmente correlacionada.

\subsection{Organização dos arquivos da segunda parte}

\paragraph{} Os arquivos empregados na segunda parte são apresentados na Tabela~\ref{tab:arquivos-parte2}.

\begin{table}[H]
    \centering
    \caption{Organização dos arquivos utilizados na identificação da planta desconhecida.}
    \label{tab:arquivos-parte2}
    \small
    \setlength{\tabcolsep}{6pt}
    \renewcommand{\arraystretch}{1.25}
    \begin{tabular}{
        @{}
        >{\raggedright\arraybackslash}p{0.25\textwidth}
        >{\raggedright\arraybackslash}p{0.68\textwidth}
        @{}
    }
        \toprule
        \textbf{Arquivo} & \textbf{Responsabilidade} \\
        \midrule
        \texttt{unkownPlant.m} & Carregar os dois conjuntos de dados, determinar a ordem e executar as duas identificações. \\
        \texttt{getOrder.m} & Calcular a solução de Wiener para ordens candidatas e gerar a curva MSE versus ordem. \\
        \texttt{getMSE.m} & Avaliar o erro médio quadrático de um vetor de coeficientes sobre todas as realizações. \\
        \texttt{adaptiveFilter.m} & Executar a recursão LMS, armazenar os erros quadráticos e estimar os coeficientes da planta. \\
        \texttt{PartTwoConsts.m} & Armazenar a ordem máxima pesquisada e os passos de adaptação. \\
        \texttt{whites.mat} & Fornecer 100 realizações com entrada branca, cada uma com 15000 amostras. \\
        \texttt{correl.mat} & Fornecer 100 realizações com entrada correlacionada, cada uma com 50000 amostras. \\
        \bottomrule
    \end{tabular}
\end{table}

\section{Modelo matemático e procedimento de identificação}

\paragraph{} Admite-se uma planta FIR causal de comprimento desconhecido \(L\), ou ordem \(M=L-1\). Para uma ordem
candidata, define-se o vetor regressor

\begin{equation}
    u(i)=
    [u(i)\;u(i-1)\;\cdots\;u(i-M)]^T.
\end{equation}

A saída observada é modelada por

\begin{equation}
    d(i)=u^T(i)w^{o}+v(i),
\end{equation}

em que \(w^{o}\) é a resposta ao impulso desconhecida e \(v(i)\) representa o ruído de
medição. Para um vetor estimado \(\widehat{w}(i)\), a saída do modelo e o erro são

\begin{equation}
    y(i)=u^T(i)\widehat{w}(i),
    \qquad
    e(i)=d(i)-y(i).
\end{equation}

O LMS atualiza os coeficientes por

\begin{equation}
    \widehat{w}(i+1)
    =
    \widehat{w}(i)
    + \mu\,u(i)e(i).
    \label{eq:lms}
\end{equation}

A curva de aprendizado é estimada pela média sobre as \(R=100\) realizações:

\begin{equation}
    \widehat{\mathrm{MSE}}(i)
    =\frac{1}{R}\sum_{r=1}^{R}e_r^2(i).
\end{equation}

\section{Determinação da ordem da planta}

\paragraph{} Para cada comprimento candidato \(L=1,2,\ldots,40\), foram estimadas a matriz de autocorrelação e a
correlação cruzada:

\begin{equation}
    \widehat{R}_{u,L}
    =\frac{1}{N_T}\sum_{r,i}u_r(i)u_r^T(i),
    \qquad
    \widehat{r}_{du,L}
    =\frac{1}{N_T}\sum_{r,i}u_r(i)d_r(i),
\end{equation}

onde \(N_T\) é o número total de amostras processadas. A solução de Wiener de cada modelo é

\begin{equation}
    \widehat{w}_{L}
    =\widehat{R}_{u,L}^{-1}\widehat{r}_{du,L}.
\end{equation}

Em seguida, o MSE de cada solução é calculado por

\begin{equation}
    \widehat{J}_{L}
    =\frac{1}{N_T}\sum_{r,i}
    \left[d_r(i)-\widehat{w}_{L}^{T}u_r(i)\right]^2.
\end{equation}

\begin{figure}[H]
    \centering
    \includegraphics[width=0.88\textwidth]{MSEordem.png}
    \caption{MSE da solução de Wiener em função da ordem candidata, utilizando o conjunto de entrada branca.}
    \label{fig:ordem-parte2}
\end{figure}

Para ordens inferiores a 30, o modelo não possui graus de liberdade suficientes para representar toda a
resposta ao impulso. O erro diminui em etapas à medida que novos coeficientes relevantes são incluídos.

A queda até aproximadamente \(-30\,\mathrm{dB}\) ocorre quando se utiliza \(L=31\). A partir desse ponto,
a inclusão de novos coeficientes não reduz o erro, que permanece limitado pelo ruído de medição. Portanto,

\begin{equation}
    M=30\quad\text{e}\quad L=31.
\end{equation}

\section{Curvas de aprendizado do MSE}

\subsection{Entrada branca}

Com \(\mu=0{,}01\), a curva média sobre as 100 realizações é mostrada na
Figura~\ref{fig:mse-branca-parte2}.

\begin{figure}[H]
    \centering
    \includegraphics[width=0.88\textwidth]{CurvaDeAprendizadoBranco.png}
    \caption{Curva de aprendizado do LMS para o caso de entrada branca.}
    \label{fig:mse-branca-parte2}
\end{figure}

A entrada branca proporciona convergência rápida. Após aproximadamente 500 iterações, a curva já se
encontra próxima do regime permanente. A média das últimas 3000 amostras é

\begin{equation}
    \mathrm{MSE}_{\infty,\mathrm{branca}}
    \approx1{,}186\times10^{-3}
    \quad\Rightarrow\quad
    -29{,}26\,\mathrm{dB}.
\end{equation}

A solução Wiener é \(9{,}9995\times10^{-4}\), praticamente \(-30\,\mathrm{dB}\). A diferença decorre do
desajuste do LMS com passo constante.

\clearpage

\subsection{Entrada correlacionada}

\paragraph{} Para a entrada correlacionada, adotou-se \(\mu=0{,}003\). A curva é apresentada na
Figura~\ref{fig:mse-corr-parte2}.

\begin{figure}[H]
    \centering
    \includegraphics[width=0.88\textwidth]{CurvaDeAprendizadoCorrelacionada.png}
    \caption{Curva de aprendizado do LMS para o caso de entrada correlacionada.}
    \label{fig:mse-corr-parte2}
\end{figure}

A convergência é consideravelmente mais lenta: a curva se aproxima do regime permanente somente após
cerca de 15000 iterações. A média das últimas 5000 amostras é

\begin{equation}
    \mathrm{MSE}_{\infty,\mathrm{corr}}
    \approx1{,}053\times10^{-3}
    \quad\Rightarrow\quad
    -29{,}77\,\mathrm{dB}.
\end{equation}

O MSE final é ligeiramente menor que no caso branco porque o passo empregado também é menor. Isso não
significa que a entrada correlacionada seja mais favorável: ela exige um número muito maior de iterações
para atingir o mesmo piso de erro.

\begin{table}[H]
    \centering
    \caption{Comparação das curvas de aprendizado.}
    \label{tab:comparacao-mse-parte2}
    \begin{tabular}{lccc}
        \toprule
        \textbf{Entrada} & \textbf{Passo \(\mu\)} & \textbf{Aproximação do regime} & \textbf{MSE em regime} \\
        \midrule
        Branca & 0,01 & cerca de 600 iterações & \(-29{,}26\,\mathrm{dB}\) \\
        Correlacionada & 0,003 & cerca de 15000 iterações & \(-29{,}77\,\mathrm{dB}\) \\
        \bottomrule
    \end{tabular}
\end{table}

\clearpage

\section{Vetor estimado da planta desconhecida}

\paragraph{} Foram consideradas três estimativas: a média dos vetores finais do LMS com entrada branca, a média dos
vetores finais com entrada correlacionada e a solução de Wiener consolidada, obtida a partir das estatísticas
dos dois conjuntos.

A distância entre as duas estimativas médias do LMS é

\begin{equation}
    \left\|
    \widehat{w}_{\mathrm{branca}}
    -\widehat{w}_{\mathrm{corr}}
    \right\|_2
    \approx1{,}61\times10^{-3},
\end{equation}

confirmando que os dois conjuntos descrevem a mesma planta. A estimativa consolidada é apresentada na
Tabela~\ref{tab:coeficientes-parte2}.

\begin{table}[H]
    \centering
    \caption{Coeficientes estimados da planta FIR de ordem 30.}
    \label{tab:coeficientes-parte2}
    \small
    \setlength{\tabcolsep}{14pt}
    \renewcommand{\arraystretch}{1.08}

    \begin{tabular}{rrr}
        \toprule
        \textbf{\(k\)}
        &
        \textbf{\(\widehat{w}_{\mathrm{white}}(k)\)}
        &
        \textbf{\(\widehat{w}_{\mathrm{corr}}(k)\)}
        \\
        \midrule
         0 &  0.4782 &  0.4783 \\
         1 &  0.1346 &  0.1364 \\
         2 &  0.2254 &  0.2259 \\
         3 & -0.1736 & -0.1738 \\
         4 & -0.0779 & -0.0778 \\
         5 & -0.0448 & -0.0444 \\
         6 &  0.1818 &  0.1811 \\
         7 & -0.0454 & -0.0457 \\
         8 &  0.1150 &  0.1150 \\
         9 & -0.3387 & -0.3359 \\
        10 & -0.0570 & -0.0580 \\
        11 & -0.1346 & -0.1356 \\
        12 & -0.2583 & -0.2596 \\
        13 &  0.0854 &  0.0826 \\
        14 &  0.0463 &  0.0469 \\
        15 &  0.0059 &  0.0053 \\
        16 & -0.2199 & -0.2177 \\
        17 &  0.1838 &  0.1852 \\
        18 &  0.0574 &  0.0581 \\
        19 & -0.0513 & -0.0504 \\
        20 &  0.0047 &  0.0044 \\
        21 & -0.0431 & -0.0430 \\
        22 & -0.2886 & -0.2867 \\
        23 & -0.0460 & -0.0451 \\
        24 & -0.1344 & -0.1359 \\
        25 & -0.1580 & -0.1603 \\
        26 & -0.1890 & -0.1894 \\
        27 & -0.0856 & -0.0879 \\
        28 & -0.3296 & -0.3287 \\
        29 &  0.1571 &  0.1570 \\
        30 &  0.0857 &  0.0850 \\
        \bottomrule
    \end{tabular}
\end{table}

\section{Conclusão}

O LMS identificou satisfatoriamente as plantas FIR utilizando apenas os pares de entrada e saída. Na primeira parte, as curvas apresentaram redução consistente do erro e valores finais compatíveis com a variância do ruído e com o erro previsto para o passo adotado.

Na segunda parte, a análise permitiu identificar uma planta FIR de ordem

\begin{equation}
    M=30,\qquad L=31.
\end{equation}

Essa ordem foi determinada pela queda do MSE da solução de Wiener até o piso de aproximadamente
\(-30\,\mathrm{dB}\). As curvas de aprendizado também confirmaram a influência das estatísticas da entrada: o sinal branco proporcionou convergência rápida, enquanto a entrada correlacionada exigiu um passo menor e aproximadamente uma ordem de grandeza a mais de iterações.

Por fim, as estimativas obtidas com os dois conjuntos de entrada apresentaram forte concordância, permitindo determinar de maneira confiável os 31 coeficientes da planta. Dessa forma, o experimento demonstrou a convergência do LMS e possibilitou a análise quantitativa do MSE e do EMSE.

\end{document}